\section{Projective modifications of abelian-surface torsors}
\label{sec:projective-modifications-torsors}

A projective modification of an algebraic space can have a scheme as
its source. In a family, however, the source need not remain flat or
have geometrically integral fibres. For torsors under abelian surfaces,
these requirements impose the following obstruction.

\begin{theorem}\label{thm:surface-torsor-modification}
Let $S$ be an arbitrary scheme, let $A$ be an abelian $S$-scheme of
relative dimension two, and let $Q$ be an fppf $A$-torsor algebraic space.
Let $X$ be a scheme, proper, flat and finitely presented over $S$, with
geometrically integral fibres. Suppose that $f:X\to Q$ is a projective $S$-morphism
and that one of the following conditions holds:
\begin{enumerate}
\item every geometric fibre $f_{\bar s}:X_{\bar s}\to Q_{\bar s}$ is
birational;
\item there is a dense open subspace $V\subset Q$ such that
$f^{-1}V\to V$ is an isomorphism.
\end{enumerate}
Then $Q$ is a scheme.
\end{theorem}

Here an abelian scheme is a smooth proper finitely presented group
scheme with geometrically connected fibres; projectivity is not
included in the definition. No group action or torsor structure
is assumed on $X$. Its fibres may be singular. The theorem therefore
includes projective modifications whose geometrically integral fibres
are geometrically normal. In the second condition,
birationality on every fibre is a consequence of the argument below;
it is not an additional hypothesis.

The proof constructs a finite locally free subscheme of $Q$ from a
rank-one pushforward. Its inclusion splits the torsor after a finite
flat base change, to which Raynaud's schematic descent criterion
applies. We give the sheaf construction over bases that may be
nonnormal and nonreduced.

The decomposition of a rank-one torsion-free sheaf into its determinant
and an ideal of points is classical. Su and Wang give its relative
Picard--Hilbert formulation for smooth spectral surfaces
\cite[Lemmas~3.22--3.23 and Remark~3.24, pp.~19--21]{SuWang},
crediting the double-dual line-bundle statement to Koll\'ar.
We give the relative sheaf argument in full, then use the resulting
finite subscheme for schematic descent of the torsor.

\begin{lemma}\label{lem:smooth-surface-flat-perfect}
Let $S$ be locally noetherian, let $Y\to S$ be a smooth algebraic space
of relative dimension two, and let $\mathcal F$ be a coherent
$\OO_Y$-module flat over $S$. Then $\mathcal F$ is perfect, locally
with a finite locally free resolution of length at most two.
\end{lemma}

\begin{proof}
Work on an affine \textup{\'etale} chart of $Y$, at a point $y$ over
$s\in S$. Put $R=\OO_{S,s}$ and let $B$ be the local ring of the chart
at $y$. The ring $B$ is flat over $R$, and $B/\mathfrak m_RB$ is a
regular local ring of dimension at most two. Write $M=\mathcal F_y$.
Choose finite free $B$-modules in exact sequences
\[
 0\longrightarrow K_1\longrightarrow P_0\longrightarrow M
 \longrightarrow0,
 \qquad
 0\longrightarrow K_2\longrightarrow P_1\longrightarrow K_1
 \longrightarrow0.
\]
The modules $K_1$ and $K_2$ are $R$-flat, and these sequences remain
exact modulo $\mathfrak m_R$. Since a regular local ring of dimension
at most two has global dimension at most two, the second syzygy
$K_2/\mathfrak m_RK_2$ is free over $B/\mathfrak m_RB$.
The fibre criterion for flatness
\cite[IV$_3$, 11.3.10]{EGA}, or \cite[Tag~00MP]{Stacks}, makes
$K_2$ flat over $B$. It is finite, hence free over the local ring $B$.
The zero-module case causes no exception. Spreading this resolution
to a neighborhood proves the assertion.
\end{proof}

\begin{lemma}\label{lem:relative-surface-rank-one-defect}
Let $S$ be locally noetherian, and let $Y\to S$ be a smooth algebraic space of relative
dimension two with geometrically integral fibres. Let $\mathcal F$ be
a coherent $\OO_Y$-module, flat over $S$, whose restriction to every
geometric fibre is torsion-free of rank one. There are an invertible
sheaf $\mathcal D$ and a closed subspace $Z\subset Y$, flat over $S$
with zero-dimensional fibres, fitting into an exact sequence
\[
 0\longrightarrow\mathcal F\longrightarrow\mathcal D
 \longrightarrow\mathcal D\otimes\OO_Z\longrightarrow0.
\]
The first map remains injective on every fibre. If $Y\to S$ is proper,
then $Z\to S$ is finite locally free.
\end{lemma}

\begin{proof}
Let $U\subset Y$ be the locus where $\mathcal F$ is locally free of
rank one. If $\mathcal F_s$ is free at a point of $Y_s$, the fibre
criterion for flatness makes $\mathcal F$ free at that point as well.
A torsion-free rank-one module over a regular local ring of dimension
zero or one is free. Hence every point of $Y_s\setminus U_s$ has local
dimension two in $Y_s$.

By Lemma~\ref{lem:smooth-surface-flat-perfect}, $\mathcal F$ has finite
locally free resolutions on \textup{\'etale} charts. The alternating
determinant of such a resolution is a line bundle with a canonical
identification with $\mathcal F$ on $U$. We explain both its gluing
and the required extension of $\mathcal F$.

At a point $y\notin U$, choose a regular sequence of length two in
the regular fibre local ring $\OO_{Y_s,y}$ and lift it to the local
ring of an \textup{\'etale} chart of $Y$. Since $Y$ is $S$-flat, the
fibrewise regular-sequence criterion
\cite[IV$_3$, 11.3.8]{EGA} makes the lifts a regular sequence in that
local ring. Thus every invertible sheaf has depth at least two there.
Similarly, choose a nonzero element of the maximal ideal of
$\OO_{Y_s,y}$. It is a nonzerodivisor on the torsion-free module
$(\mathcal F_s)_y$. Flatness of $\mathcal F$ over $S$ and the same
criterion show that a lift is a nonzerodivisor on $\mathcal F_y$.
Consequently $\mathcal F_y$ has depth at least one.

Put $j:U\hookrightarrow Y$. The depth-one assertion means that no
associated point of $\mathcal F$ lies outside $U$, so
\[
 \mathcal F\longrightarrow j_*(\mathcal F|_U)
\]
is injective. Indeed, a nonzero kernel would have an associated point
of $\mathcal F$ in the complement. The depth-two extension lemma
\cite[Tag~0E9I]{Stacks} says that any invertible sheaf $\mathcal L$
on a chart satisfies $\mathcal L=j_*(\mathcal L|_U)$ there.
This lemma requires no normality or reducedness.

On a double overlap of two charts, the determinant lines have an
isomorphism over $U$ through their identifications with $\mathcal F$.
Their Hom sheaf is invertible. The depth-two extension lemma extends
that isomorphism uniquely; it also extends its inverse. Restriction
is injective, so the extended maps are inverse and satisfy the cocycle
identity on triple overlaps. \textup{\'Etale} descent therefore gives
an invertible sheaf $\mathcal D$ on $Y$ with
$\mathcal D|_U=\mathcal F|_U$. This construction does not require a
global locally free resolution on $Y$. The same extension lemma gives
\[
 \mathcal F\lhook\joinrel\longrightarrow
 j_*(\mathcal F|_U)=j_*(\mathcal D|_U)=\mathcal D.
\]

For every geometric $s$, this map restricts to an isomorphism on the
dense open $U_s$ of the integral surface $Y_s$. Its fibre kernel has
rank zero, so torsion-freeness of $\mathcal F_s$ makes that kernel
zero. Let $\mathcal C$ be the cokernel on $Y$. Since $\mathcal D$ is
$S$-flat, the fibrewise injectivity gives
\[
 \operatorname{Tor}_1^{\OO_{S,s}}
       (\mathcal C_y,\kappa(s))=0
 \quad\text{for every }y\text{ over }s.
\]
The local criterion for flatness \cite[Tag~00MK]{Stacks} applies to
the finite module $\mathcal C_y$ over the noetherian local ring at
$y$, and shows that $\mathcal C$ is $S$-flat.

Tensoring the inclusion by $\mathcal D^{-1}$ gives an ideal
$\mathcal I\subset\OO_Y$. Define $Z$ by $\mathcal I$. Then
$\mathcal C=\mathcal D\otimes\OO_Z$, so $Z$ is flat over $S$.
Its fibres are supported in $Y_s\setminus U_s$, which has dimension
zero. If $Y/S$ is proper, $Z/S$ is proper and quasi-finite, hence
finite \cite[Tag~0A4X]{Stacks}. It is then finite locally free.
\end{proof}

\begin{lemma}\label{lem:ample-pushforward-flat-basechange}
Let $S$ be noetherian affine, let $Y$ be a quasi-compact separated
algebraic space of finite type over $S$, and let $f:X\to Y$ be proper
with an $f$-ample invertible sheaf $L$. Suppose that $X$ is a scheme,
flat and of finite type over $S$. For sufficiently large $n$, the
coherent sheaf
\[
 \mathcal F=f_*L^n
\]
is flat over $S$, its formation commutes with arbitrary base change
on $S$, and the evaluation $f^*\mathcal F\to L^n$ is surjective.
\end{lemma}

\begin{proof}
Relative Serre vanishing and generation, applied on a finite affine
\textup{\'etale} cover of $Y$, give a common $n$ for which
$R^if_*L^n=0$ for $i>0$ and the evaluation map is surjective.
Properness makes $\mathcal F$ coherent.

Let $V\to Y$ be an affine \textup{\'etale} chart. The scheme $X_V$
is quasi-compact, separated and $S$-flat. A finite affine cover of
$X_V$ has affine finite intersections and gives a bounded ordered
\v Cech complex $C^\bullet$ computing the cohomology of $L^n$ on
$X_V$. Every term is flat over $\OO(S)$: the coordinate ring of each
affine intersection is flat over $\OO(S)$, and the module associated
to the restricted line bundle is finite projective over that ring.
The positive cohomology of $C^\bullet$ vanishes by the chosen
relative Serre vanishing and flat base change to $V$.

Starting at the final nonzero term, exactness and flatness of the
successive quotients show that every successive kernel is flat,
including $H^0(C^\bullet)$. The corresponding short exact sequences
remain exact after tensoring with any $\OO(S)$-module. Thus tensoring
the complex preserves its vanishing in positive degrees and computes
the tensor product of $H^0(C^\bullet)$ in degree zero.
After any affine base change $S'\to S$, the base-changed affine cover
has exactly this tensor-product \v Cech complex. This proves both
flatness and the base-change isomorphism
\[
 \mathcal F|_{Y_{S'}}\simeq
 (f_{S'})_*\bigl(L^n|_{X_{S'}}\bigr).
\]
The assertions glue across the \textup{\'etale} cover of $Y$ and
across an affine cover of $S'$. In particular $S'\to S$ need not be
flat or noetherian.
\end{proof}

\begin{proof}[Proof of Theorem~\ref{thm:surface-torsor-modification}]
We first assume that $S$ is noetherian affine and condition~(1) holds.
Choose an $f$-ample line bundle $L$ and take $n$ as in
Lemma~\ref{lem:ample-pushforward-flat-basechange}. Write
$\mathcal F=f_*L^n$. On every geometric fibre,
\[
 \mathcal F_{\bar s}=(f_{\bar s})_*L_{\bar s}^n
\]
is torsion-free of rank one. Birationality gives rank one. To check
torsion-freeness, a section killed by a nonzero function on the
integral target vanishes over a dense open. Its pullback is a section
of a line bundle on the integral $X_{\bar s}$, so it is zero.

Apply Lemma~\ref{lem:relative-surface-rank-one-defect} to obtain an
invertible $\mathcal D$ and a finite locally free closed subscheme
$Z\subset Q$ with $\mathcal F=\mathcal D\otimes\mathcal I_Z$.
For every geometric $s$, we claim
\[
 Z_{\bar s}=\varnothing
 \quad\Longleftrightarrow\quad
 f_{\bar s}\text{ is an isomorphism}.
\]
If $f_{\bar s}$ is an isomorphism then $\mathcal F_{\bar s}$ is
invertible. The fibre criterion makes $\mathcal F$ invertible at
all points of that fibre, so $Z_{\bar s}$ is empty.
Conversely, if $Z_{\bar s}$ is empty, the evaluation map is a
surjection of line bundles
\[
 f_{\bar s}^*\mathcal D_{\bar s}\longrightarrow L_{\bar s}^n,
\]
hence an isomorphism. The line bundle $L_{\bar s}^n$ is therefore
trivial on each fibre of $f_{\bar s}$. Since it is relatively ample,
these proper fibres are zero-dimensional. Thus $f_{\bar s}$ is
finite. It is birational onto the normal surface $Q_{\bar s}$ and
has integral source, so it is an isomorphism.

The rank of $Z/S$ is locally constant. Decompose $S$ into the open
and closed parts $S_0$ and $S_+$ on which that rank is respectively
zero and positive. Over $S_0$ every fibre map is an isomorphism, and
so is $f$. Indeed, its non-quasi-finite locus is empty, making $f$
finite. The coherent cokernel of $\OO_Q\to f_*\OO_X$ has zero fibres,
so it is zero by Nakayama's lemma. The source and target of this
surjection are $S$-flat; hence its kernel also has zero fibres and
is zero by Nakayama's lemma. It follows that $Q_{S_0}\simeq X_{S_0}$
is a scheme. This also explains why, over a local base, an isomorphism
on the closed fibre cannot hide a nonisomorphic fibre elsewhere:
the proper image of the closed non-quasi-finite locus would avoid
the closed point, so that locus would be empty. For the resulting
finite map, the cokernel of $\OO_Q\to f_*\OO_X$ vanishes along the
closed fibre by Nakayama. Its support is proper over the local base,
and therefore it is empty. The kernel then vanishes by the same
argument, using $S$-flatness to preserve exactness on the closed fibre.

Over $S_+$, the map $Z\to S_+$ is finite locally free and surjective.
Its inclusion in $Q$ gives a section of $Q_Z$, so $Q_Z\simeq A_Z$.
An abelian scheme over an affine base is AF, meaning that every finite
set of its points lies in an affine open; see the application following
\cite[Corollary~3.8, p.~5-15]{RaynaudDescent}. Finite morphisms preserve AF on
their source. Thus, after restricting to an affine part of $S_+$,
$A_Z$ is AF because it is finite over $A$.

Apply \cite[Theorem~4.2, pp.~5-16--5-18]{RaynaudDescent} with
$Z\to S_+$ as both its fpqc cover and its finite cover. On that cover
$Q_Z=A_Z$ is a noetherian scheme, and its finite descent orbits have
affine neighborhoods because it is AF. Therefore $Q_{S_+}$ is a
scheme. This proves the noetherian case under condition~(1).

We now handle both conditions over an arbitrary base. Work on an
affine open $S=\Spec R$ and write $R$ as a filtered union of finitely
generated $\mathbf Z$-subalgebras. Descend $A$ and $X$ as schemes,
$Q$ as an algebraic space, and the morphism $f$, the group law,
the action, and the torsor isomorphism, together with its inverse.
The descent of finitely presented algebraic spaces and their maps
is \cite[Tag~07SK]{Stacks}; the scheme version is
\cite[IV$_3$, 8.8--8.10]{EGA}. All finitely many identities hold
after increasing the index.

One can retain the torsor property by also descending an actual
trivializing cover. Choose an affine scheme $W$ with an \textup{\'etale}
surjection $W\to Q$ of finite presentation. Then $W\to S$ is smooth
and surjective, and $W\to Q$ gives a section of $Q_W$. Descend $W$,
this section, and the resulting isomorphism $Q_W\simeq A_W$ with
its inverse. Retain smoothness and surjectivity of $W_i\to S_i$.
The descended $Q_i$ is consequently an $A_i$-torsor; in particular
it is proper and smooth over $S_i$.

Retain properness of the scheme morphisms $X_i\to S_i$ and
$A_i\to S_i$ by \cite[IV$_3$, 8.10.5]{EGA}, flatness by
\cite[IV$_3$, 11.2.6(ii)]{EGA}, and smoothness of $A_i$ and $W_i$
by \cite[IV$_4$, 17.7.8(ii)]{EGA}. Geometric integrality of the
$X_i$-fibres and geometric connectedness of the $A_i$-fibres can
also be retained. The bad loci on a noetherian model are constructible
by \cite[IV$_3$, 9.7.7]{EGA}, and their inverse images in $S$ are
empty. Each bad locus is a finite union of locally closed pieces
$\Spec((R_i/I)[1/a])$. The tensor product of the coordinate ring
of such a piece with $R$ is zero; the equation $1=0$ already holds
at a finite stage. Increasing the index removes all the pieces.
The relative dimension of $A_i$ can be fixed to two by removing its
other open-and-closed dimension loci, whose pullback to $S$ is empty.

Finally descend $L$ as an invertible sheaf and retain its relative
ampleness. Choose a finite affine \textup{\'etale} cover of a model
of $Q$ and pull it back through the system. On each such affine
target, relative ampleness is ampleness on the inverse-image scheme.
The limit criterion \cite[IV$_3$, 8.10.5.2]{EGA}, equivalently
\cite[Tag~09MT]{Stacks}, makes $L_i$ ample there at a common finite
stage. Relative ampleness is \textup{\'etale} local on the target,
so $L_i$ is $f_i$-ample. The map $f_i$ is proper because $X_i$ is
proper over $S_i$ and $Q_i$ is separated over $S_i$. At this point
no birationality assertion about $f_i$ has been imposed.

On this noetherian model choose $n$ by
Lemma~\ref{lem:ample-pushforward-flat-basechange} and put
\[
 \mathcal F_i=(f_i)_*L_i^n.
\]
This sheaf is $S_i$-flat and commutes with arbitrary base change.
By Lemma~\ref{lem:smooth-surface-flat-perfect} it is perfect on
$Q_i$. Its virtual rank $r_i$, the alternating sum of the ranks in
a finite locally free resolution, is locally constant and is preserved
by base change. Since the fibres of $Q_i/S_i$ are geometrically
connected, $r_i$ is constant on each fibre. It therefore comes from
a locally constant function $\rho_i$ on $S_i$: the images of its
disjoint rank loci are disjoint opens, because $Q_i\to S_i$ is
smooth, open and surjective.

Under condition~(1), the base change $\mathcal F=f_*L^n$ has generic
rank one on every geometric fibre, hence virtual rank one everywhere.
To see this directly, a finite locally free resolution remains exact
after base change on $S_i$, since $\mathcal F_i$ and its syzygies
are $S_i$-flat. At the generic point of an integral fibre its Euler
rank is the dimension of its sole cohomology module. Under
condition~(2), $\mathcal F$ is invertible on the dense open $V$ where
$f$ is an isomorphism. Its locally constant virtual rank is then one
on a dense subset, hence everywhere. In either case $S\to S_i$
factors through the open-and-closed locus where $\rho_i=1$.
Replace $S_i$ by that locus.

Let $\bar s$ be a geometric point of the resulting model, and let
$\eta$ be the generic point of the integral surface $Q_{i,\bar s}$.
Flatness and base change give
\[
 \dim_{\kappa(\eta)}
 H^0\bigl((X_i)_\eta,L_i^n|_{(X_i)_\eta}\bigr)=1.
\]
Indeed, base-changing a finite locally free resolution to $\bar s$
and then localizing at $\eta$ gives a complex of finite-dimensional
vector spaces with sole cohomology $(\mathcal F_i)_{\bar s,\eta}$,
so that dimension equals its virtual rank one.
The evaluation map is surjective; consequently the single nonzero
global section generates the line bundle on $(X_i)_\eta$, making
it trivial. It is also ample. A proper scheme over a field with
ample trivial line bundle is zero-dimensional and hence finite.
Its function algebra has dimension one, so
\[
 (X_i)_\eta=\Spec\kappa(\eta).
\]
The integral source $X_{i,\bar s}$ therefore maps birationally to
$Q_{i,\bar s}$. This holds for every geometric $\bar s$ of the
model. The noetherian case already proved applies, making $Q_i$ a
scheme. Its base change $Q$ is a scheme. The affine opens of $S$
cover the original base, which proves the theorem.
\end{proof}

\begin{remark}\label{rem:surface-modification-scope}
The theorem excludes projective scheme modifications of nonschematic
abelian-surface torsors that retain flatness and geometrically integral
fibres. It does not assert schematicity of the full Picard functor of
an arbitrary proper flat family. The surface hypothesis enters when
the rank-one defect has zero-dimensional fibres. In higher relative
dimension the defect can be positive-dimensional, so the argument
does not supply the finite splitting cover required by schematic
descent. The proof also shows that projectivity of $f$ may be replaced
by properness together with an $f$-ample invertible sheaf.
\end{remark}
